\section{Experiments}
\label{sec:experiments}

We evaluate MI on two tasks chosen to probe different capabilities: reactive control (CartPole) and static classification (MNIST). The goal is not to claim state-of-the-art accuracy on any single task but to demonstrate that a self-organizing developmental architecture can solve non-trivial problems and to characterize its distinctive behaviors.

\subsection{CartPole-v1}
\label{sec:exp_cartpole}

\paragraph{Setup.}
OpenAI Gym CartPole-v1, 4 observations (cart position, velocity, pole angle, angular velocity) encoded into 32 sensory morphons via Gaussian population coding (8 receptive-field tiles per observation). 2 motor morphons (left/right force). Cerebellar developmental preset: seed\_size=60, dimensions=4, initial\_connectivity=0.25.

A small linear TD-error critic external to the MI network provides the reward modulation signal, biologically modelling dopamine neurons in VTA/SNc that compute $\delta = R + \gamma V(s') - V(s)$ from striatal input \cite{fremaux2013reinforcement}. The critic is the only component using exact gradients; everything inside the MI network is trained by three-factor learning.

\paragraph{Result.}
\textbf{Solved} on both v0.5.0 and v4.1.0 standard profiles. v0.5.0: avg(100)=195.2 at episode 895, best=468. v4.1.0: avg(100)=195.5 at episode 936, best=468. The Pulse Kernel Lite and iSTDP changes introduced for MNIST did not affect CartPole learning dynamics. The solve confirms that the full MI pipeline --- developmental bootstrap, three-factor learning, tag-and-capture consolidation, contrastive reward, Endoquilibrium regulation --- is stable across major architectural changes.

\begin{figure}[h]
  \centering
  \figifexists[0.85\linewidth]{cartpole_curve.pdf}{CartPole learning curve}
  \caption{CartPole-v1 learning curve (v4.1.0 standard profile, 1000 episodes). SOLVED at episode 936, avg(100)=195.5. The solve is consistent across versions: v0.5.0 solved at ep 895 (avg=195.2); v4.1.0 solves at ep 936 (avg=195.5).}
  \label{fig:cartpole}
\end{figure}

\paragraph{What this proves.}
The full MI pipeline --- developmental bootstrap, three-factor learning, tag-and-capture consolidation, contrastive reward, Endoquilibrium regulation --- works end-to-end on a non-trivial control task. CartPole is the prerequisite gate for all other benchmarks: if the basic learning loop is broken, nothing else can work. The CartPole solve provides the foundation to investigate harder tasks.

\subsection{MNIST}
\label{sec:exp_mnist}

\paragraph{Setup.}
Standard MNIST, 784 input pixels (28$\times$28), 10 output classes. Cortical developmental preset: seed\_size=200, dimensions=6, initial\_connectivity=0.02 (Phase 4 hardcodes 30\% S$\to$A density). Three training epochs of 5000 samples each (standard profile) or 3000 (quick profile). Analog readout (Section~\ref{sec:implementation}). Endoquilibrium enabled with the corrected stage-detection thresholds described in Section~\ref{sec:failure_modes}. Inhibitory competition via LocalInhibition mode with Vogels 2011 iSTDP (Section~\ref{sec:istdp}); fast path via Pulse Kernel Lite (Section~\ref{sec:pulse_kernel}). The training pattern follows Diehl \& Cook \cite{diehl2015unsupervised}: rate-coded Poisson input, unsupervised STDP-driven feature formation, and post-hoc neuron labeling for evaluation. All results from v4.1.0, seed=42.

\paragraph{Analog readout rationale.}
A key design choice is using an analog readout (a linear layer over morphon potentials) rather than reading classification decisions from spike counts at motor morphons. The motivation: the spike pipeline destroys discriminative information. Morphon potentials carry the network's learned representation, but conversion to spikes, propagation through synaptic delays, leaky integration, and multi-hop accumulation at motor morphons each introduce lossy compression. A simple synthetic 2-class probe (bag-of-characters encoding, 27-dim input) confirmed this empirically: a spike-based motor readout on the same trained network produces chance-level accuracy ($\sim$46\%) while a linear analog readout over the same morphon potentials reaches 99\%. The information that distinguishes classes is present in the substrate --- the analog readout proves it --- but the spike pipeline cannot deliver it to motor morphons intact. The analog readout is therefore not a shortcut; it is the correct output mechanism for a spiking architecture where the spike pathway serves feature formation, not classification reporting. Biology achieves the same separation: Purkinje cells in the cerebellum perform analog integration of climbing and parallel fiber inputs rather than thresholding to binary spikes for motor control \cite{ito2006cerebellar}.

\paragraph{Results.}
\begin{table}[h]
\centering
\caption{MNIST results across versions and training conditions ($\sim$1240--1290 morphons, $\sim$73--88K synapses). v4.1.0, seed=42. Accuracy varies $\pm$2--3pp across independent runs due to stochastic synaptogenesis.}
\label{tab:mnist}
\begin{tabular}{lccc}
\toprule
Condition & Quick (3K×3ep) & Standard (5K×3ep) & Notes \\
\midrule
V2 GlobalKWTA (intact)    & 44--44.5\% & 48.5\% & Algorithmic k-WTA \\
V3 LocalInhib+iSTDP (intact) & 44.5--48\% & \textbf{52.0\%} & Vogels 2011; $\geq$V2; 4.3×--4.4× faster \\
Post-damage (30\% killed)    & 48.0\%     & 47.3\%  & Random ablation \\
Post-recovery (1500 samples) & 33.5\%     & 28.5\%  & See §~\ref{sec:self_healing_discussion} \\
\midrule
\textit{v3.0.0 V2 (rate iSTDP, intact)} & \multicolumn{2}{c}{\textit{30.0\%}} & \textit{Hub-dominated; reference baseline} \\
\textit{v3.0.0 post-recovery}            & \multicolumn{2}{c}{\textit{48.0\%}}  & \textit{+18pp; hub-clearing via damage} \\
\bottomrule
\end{tabular}
\end{table}

\paragraph{Hub dominance and its cure.}
The v3.0.0 baseline (30\%) was limited by a specific failure mode: a small number of ``hub'' morphons with very high in-degree accumulated nearly all reward-correlated weight updates and began responding to every input class, collapsing the readout to a near-random discriminator. With rate-based inhibitory STDP --- computed once per 10-tick medium period using a 100-tick activity average --- the inhibitory interneurons responded too slowly to break this dominance; a single hub could accumulate 50--100 correlated spikes before any corrective signal arrived.

Replacing rate-based iSTDP with the spike-timing rule of \citet{vogels2011inhibitory} fixes this at the root. On each delivered inhibitory spike (pre-spike): $\Delta w = -\eta \cdot x_{\text{post}}$, where $x_{\text{post}}$ is the post-synaptic trace. On each excitatory post-spike: $\Delta w = -\eta \cdot (x_{\text{pre}} - \alpha)$, where $\alpha = 2\rho_0\tau$ sets the equilibrium firing rate. The result: a hub morphon that fires at $5\times$ the target rate immediately recruits stronger inhibition on the very spikes that evidence its over-activity, rather than waiting for the slow running average to catch up. The readout response after this fix: 1 labeled class (class 4 only) $\rightarrow$ 12 labeled classes across all 10 digits, and accuracy 8\% $\rightarrow$ 44.5\%.

\paragraph{The self-healing finding revisited.}
\label{sec:self_healing_discussion}

In v3.0.0 (rate-based iSTDP), post-damage recovery reached \textbf{48\%} --- 18pp above the 30\% intact baseline. The mechanism: damage randomly removed entrenched hub morphons, Endoquilibrium responded by re-entering the high-plasticity Differentiating stage, and fresh morphons grew into well-specialized representations the hubs had blocked. The damage-and-regrow cycle was a manual proxy for hub clearing.

In v4.1.0, spike-timing iSTDP prevents hub formation during normal training, raising the intact ceiling to \textbf{52\%} on the standard profile. Post-damage recovery (1500 samples) reaches 28.5\% --- well below the intact baseline. This is expected: with no entrenched hubs to clear, damage is pure loss. The self-healing effect relied entirely on the pathology it was fixing.

This reframes the finding. The v3.0.0 self-healing result was not a feature of the architecture; it was a diagnostic. The damage cycle revealed that hub morphons were the limiting factor and that the substrate was capable of 48\% under the right conditions. That diagnostic pointed directly to the fix: inhibitory plasticity operating at the right timescale. With the fix in place, the system reaches 52\% without damage. The +18pp self-healing gain has been absorbed into the intact baseline.

\begin{figure}[h]
  \centering
  \figifexists[0.7\linewidth]{self_healing.pdf}{self-healing bar chart}
  \caption{MNIST across versions. v3.0.0 (rate iSTDP): 30\% intact, 48\% post-recovery --- damage cleared hubs that training could not. v4.1.0 (spike-timing iSTDP): 52\% intact, 28.5\% post-recovery --- iSTDP prevents hub formation, making the damage trick unnecessary. The +18pp self-healing gain has been absorbed into the intact baseline.}
  \label{fig:self_healing}
\end{figure}

\paragraph{Implication.}
The intact ceiling has risen from 30\% to 52\% (standard profile) via a single architectural correction: inhibitory plasticity at biological timescales. No tuning was required. The 22pp gain is entirely attributable to the iSTDP timescale fix; everything else (network size, training data, epochs) is unchanged.

\begin{figure}[h]
  \centering
  \figifexists[0.7\linewidth]{receptive_fields.pdf}{MNIST receptive fields}
  \caption{Receptive fields of the top-12 associative morphons (by mean incoming weight magnitude) after MNIST training, dumped from the v3.0.0 quick-profile run. Each tile shows the morphon's S$\to$A synapse weights mapped to the 28$\times$28 pixel grid (red = positive, blue = negative). Many morphons develop spatially-localized, digit-stroke-like receptive fields without any backprop-driven gradient signal --- features emerge from STDP + k-WTA + reward-modulated three-factor learning.}
  \label{fig:receptive_fields}
\end{figure}
